\documentclass[journal,twoside]{IEEEtran}

\usepackage{amsmath,amssymb}
\usepackage{graphicx}
\usepackage{booktabs}
\usepackage{multirow}
\usepackage{cite}

\begin{document}

\title{Semantic Affinity Guided Semi-Supervised Surgical Scene Segmentation}

\author{Anonymous Authors}

\maketitle

\begin{abstract}
This document describes the proposed method in an IEEE double-column journal format. The method extends UniMatch-V2 for surgical semantic segmentation by introducing a dataset-specific vision-language affinity prior, adaptive logit fusion, class-imbalance debiasing, and boundary-aware geometric regularization. The central motivation is that surgical datasets contain long-tailed, visually similar, and spatially small categories, for which pixel-level supervision alone often provides insufficient semantic discrimination.
\end{abstract}

\begin{IEEEkeywords}
Semi-supervised segmentation, surgical scene understanding, vision-language alignment, long-tailed learning, boundary regularization.
\end{IEEEkeywords}

\section{Method}

\subsection{Overview}
Let $\mathcal{D}_l=\{(x_i,y_i)\}_{i=1}^{N_l}$ denote the labeled set and $\mathcal{D}_u=\{u_i\}_{i=1}^{N_u}$ denote the unlabeled set, where $y_i \in \{1,\ldots,C\}^{H\times W}$ is a pixel-wise surgical semantic mask. Our framework is built upon UniMatch-V2, but introduces a semantic affinity branch to compensate for rare-class under-learning and small-object ambiguity. The method contains two stages. First, we learn a dataset-specific patch-to-text affinity model using DINOv2 visual tokens and SigLIP textual class anchors. Second, the learned affinity model is attached to the semi-supervised segmentation network as a lightweight side branch and fused with dense segmentation logits by a learnable spatial gate.

The key design is to decouple image-level semantic alignment from pixel-level semi-supervised segmentation while keeping both stages in a shared DINOv2 feature space. This avoids training an independent visual backbone for the auxiliary prior and enables the prior to be injected into the segmentation model with minimal inference overhead.

\subsection{Stage I: Dataset-Specific Vision-Language Affinity Prior}
For each surgical dataset, we construct a class vocabulary $\mathcal{C}=\{c_1,\ldots,c_C\}$. Each class name is expanded into several medical text templates, and the SigLIP text encoder maps these prompts into semantic embeddings. After template ensembling and normalization, a class anchor $t_c$ is obtained for class $c$.

Given an image $x$, the DINOv2 encoder extracts patch tokens
\begin{equation}
    P = f_{\theta}^{\mathrm{vis}}(x) \in \mathbb{R}^{N \times d},
\end{equation}
where $N$ is the number of image patches. Patch tokens and text anchors are projected into a shared latent space by lightweight projection heads:
\begin{equation}
    \hat{p}_i = \frac{g_v(p_i)}{\|g_v(p_i)\|_2}, \quad
    \hat{t}_c = \frac{g_t(t_c)}{\|g_t(t_c)\|_2}.
\end{equation}
The patch-class affinity is then computed as
\begin{equation}
    A_{i,c} = \frac{\hat{p}_i^\top \hat{t}_c}{\tau},
\end{equation}
where $\tau$ is a learnable temperature. To obtain image-level class evidence, we aggregate the top-$k$ patch affinities:
\begin{equation}
    s_c = \frac{1}{k}\sum_{i \in \mathrm{TopK}(A_{:,c})} A_{i,c} + b_c ,
\end{equation}
where $b_c$ is a learnable class bias. The affinity model is trained with multi-label binary cross entropy using image-level labels derived from segmentation masks:
\begin{equation}
    \mathcal{L}_{\mathrm{aff}} =
    -\sum_{c=1}^{C} w_c
    \left[z_c \log \sigma(s_c) + (1-z_c)\log(1-\sigma(s_c))\right],
\end{equation}
where $z_c$ indicates whether class $c$ is present in the image and $w_c$ compensates class-frequency imbalance.

Unlike a generic vision-language prior, the proposed affinity prior is trained separately for each dataset. This is important because surgical datasets differ substantially in instrument taxonomy, organ appearance, annotation granularity, and rare-class distribution. The resulting checkpoint stores the visual projection head, text projection head, normalized class anchors, temperature, and class bias.

\subsection{Stage II: Affinity-Guided Semi-Supervised Segmentation}
In the second stage, we train a UniMatch-V2 segmentation model with a DPT decoder over DINOv2 features. For labeled data, the supervised objective is the standard pixel-wise cross entropy:
\begin{equation}
    \mathcal{L}_{x} = \mathrm{CE}(F_{\phi}(x), y),
\end{equation}
where $F_{\phi}$ denotes the segmentation network.

For unlabeled images, an exponential moving average teacher predicts weak-view pseudo labels. Let $q^w$ denote teacher logits on a weakly augmented image and $\hat{y}=\arg\max \mathrm{softmax}(q^w)$ be the pseudo label. Two strongly augmented views are optimized by consistency learning:
\begin{equation}
    \mathcal{L}_{u} =
    \frac{1}{2}\sum_{m=1}^{2}
    \mathbf{1}[\max q^w \ge \eta]\,
    \mathrm{CE}(F_{\phi}(u_m^s), \hat{y}),
\end{equation}
where $\eta$ is the confidence threshold.

\subsection{Semantic Affinity Logit Fusion}
The proposed side branch reuses the DINOv2 patch tokens from the segmentation model. Given patch tokens $P$, the frozen affinity module produces a dense prior map $R \in \mathbb{R}^{C \times H \times W}$ by computing patch-class affinities and upsampling them to the image resolution. A learnable $1\times1$ convolution predicts a pixel-wise gate:
\begin{equation}
    G = \sigma(h(P)) \in [0,1]^{C \times H \times W}.
\end{equation}
The final segmentation logits are
\begin{equation}
    L_{\mathrm{final}} =
    (1-G \odot M) \odot L_{\mathrm{seg}}
    + (G \odot M) \odot R,
\end{equation}
where $L_{\mathrm{seg}}$ is the DPT segmentation logit, $M$ is a non-background mask that prevents the prior from dominating the background channel, and $\odot$ denotes element-wise multiplication. The gate bias is initialized to a negative value so that the network initially relies on the segmentation branch and gradually learns when to trust the semantic prior.

This fusion is the main innovation of the framework. Rather than using text embeddings only for image-level auxiliary classification, the method converts patch-text similarity into a dense spatial prior and injects it directly into semi-supervised segmentation. This is especially beneficial for rare instruments, small anatomical structures, and visually ambiguous categories whose pseudo labels are often suppressed by majority classes.

\subsection{Class Prior Debiasing}
Surgical datasets are highly long-tailed. To reduce pseudo-label collapse toward frequent classes, we estimate the class prior $\pi_c$ from labeled masks and adjust teacher logits before pseudo-label selection:
\begin{equation}
    \tilde{q}_c = q_c - \lambda \log \pi_c ,
\end{equation}
where $\lambda$ is warmed up during early epochs. This logit adjustment raises the relative score of rare classes and improves the probability that their pixels survive confidence thresholding.

\subsection{Boundary and Tangent-Space Geometry Regularization}
Pixel-wise pseudo labels often produce coarse object boundaries, which is problematic for surgical instruments and thin anatomical structures. We therefore add a boundary auxiliary head that predicts a one-channel boundary logit from segmentation logits. The target boundary map is generated from ground-truth or pseudo masks using Sobel operators. The boundary loss combines binary cross entropy and Dice loss:
\begin{equation}
    \mathcal{L}_{b} =
    (1-\alpha_b)\mathcal{L}_{\mathrm{BCE}}
    + \alpha_b \mathcal{L}_{\mathrm{Dice}}.
\end{equation}
For CutMix augmented unlabeled samples, artificial paste outlines are excluded from the boundary loss so that the model does not learn augmentation seams as anatomical boundaries.

In addition, we apply a tangent-space consistency regularizer to align local boundary geometry between predictions and targets. This encourages coherent boundary orientation and suppresses fragmented contours. The geometry term is denoted as $\mathcal{L}_{t}$.

\subsection{Multi-Label Classification Auxiliary Loss}
Because surgical images usually contain multiple instruments or anatomical regions, we further attach a multi-label classification head to the labeled stream. The target vector is derived from mask-level class presence. We use asymmetric loss to reduce the dominance of easy negative labels:
\begin{equation}
    \mathcal{L}_{c} = \mathrm{ASL}(r, z),
\end{equation}
where $r$ is the image-level prediction and $z$ is the multi-hot label. The classification loss is warmed up to avoid destabilizing early dense prediction.

\subsection{Overall Objective}
The final training objective is
\begin{equation}
    \mathcal{L} =
    \frac{1}{2}(\mathcal{L}_{x}+\mathcal{L}_{u})
    + \beta_b \mathcal{L}_{b}
    + \beta_t \mathcal{L}_{t}
    + \beta_c \mathcal{L}_{c}
    + \beta_a \mathcal{L}_{\mathrm{aff}}^{\mathrm{aux}},
\end{equation}
where $\mathcal{L}_{\mathrm{aff}}^{\mathrm{aux}}$ is the image-level BCE loss computed from the affinity side branch on labeled images. The affinity projection and anchors are frozen at the beginning of Stage II. When validation mIoU saturates, the projection head can be unfrozen with a small learning rate, allowing the prior to adapt to the segmentation feature distribution while preserving the learned semantic class geometry.

\subsection{Inference}
At inference time, each image is passed through DINOv2 once. The same patch tokens feed both the DPT decoder and the affinity side branch. The final mask is obtained by applying argmax to the gated fused logits:
\begin{equation}
    \hat{y} = \arg\max_c L_{\mathrm{final},c}.
\end{equation}
The additional cost of the proposed branch is limited to patch-anchor matrix multiplication and a $1\times1$ gating convolution, making the method suitable for surgical segmentation settings where computational efficiency is required.

\section{Innovation Summary}
The proposed method differs from conventional semi-supervised segmentation in four aspects. First, it learns a dataset-specific vision-language affinity prior instead of using a generic frozen language model. Second, it transforms image-level patch-text alignment into dense segmentation priors. Third, it injects the prior through a conservative learnable gate rather than hard pseudo-label replacement. Fourth, it combines semantic prior fusion with long-tail logit debiasing and boundary-tangent geometry regularization, addressing class imbalance, semantic ambiguity, and boundary imprecision in a unified framework.

\end{document}
